% Document settings
\documentclass[11pt]{book}

% Key Organizational Parts in the Book Format
% \frontmatter:
% This command is used to generate the initial section of the book that typically contains the title page, dedication, preface, and acknowledgments. It formats the page numbering in Roman numerals.
% \mainmatter:
% This command marks the beginning of the main content of the book. Page numbering switches to Arabic numerals here. This is where your chapters and sections are defined.
% \backmatter:
% This command indicates the concluding section of the book. It usually includes the index, glossary, bibliography, and any appendices. Page numbering will continue in Arabic numerals.
% \chapter:
% This command creates a new chapter. Chapters are the highest-level division in the book class and start on a new page.
% \section:
% This command creates a new section within a chapter. Sections help organize content within a chapter.
% \subsection:
% This command creates a subsection within a section, further dividing the content for clarity.
% \subsubsection:
% This command creates a subsubsection within a subsection for even finer divisions of content.
% \paragraph:
% This command creates a paragraph heading, which is a smaller division within a section or subsection. It does not start a new line by default.
% \subparagraph:
% This creates an even smaller division within a paragraph. Like \paragraph, it does not start a new line by default.

\usepackage[letterpaper,margin=.75in,inner=.75in,outer=2in]{geometry}
%\usepackage{pdflscape}
\usepackage{graphicx}
\usepackage{subcaption}
\usepackage{wrapfig}
\usepackage{multirow}
\usepackage{setspace}
\usepackage[bookmarks=true, colorlinks=true, linkcolor=black, citecolor=black, urlcolor=black]{hyperref}
\usepackage{xurl}
\usepackage{booktabs}
\usepackage{tcolorbox}
\usepackage{listings}          % for creating language style
\usepackage{xcolor} % For coloring
\usepackage{textcomp} % For additional symbols
\usepackage{marginnote}
\pagestyle{plain}

\setlength{\parskip}{10pt}
\setlength\parindent{0pt}

%\usepackage{minitoc}

%\usepackage{boisik}


% Prevent lstlisting from breaking between pages
\lstset{
    breaklines=true,                  % Allow line breaks within long lines
    keepspaces=true,                  % Keep spaces in listings
    prebreak=\mbox{$\hookleftarrow$},   % Add a marker for the line breakpoint, carriagereturn
    escapeinside={(*@}{@*)},          % Control escape inside code
}

\input{ArduinoLatexListing/arduinoLanguage.tex}    % adds the arduino language listing

%% Define an Arduino style fore use later %%
\lstdefinestyle{myArduino}{
  language=Arduino,
  %% Add other words needing highlighting below %%
  morekeywords=[1]{},                  % [1] -> dark green
  morekeywords=[2]{FILE_WRITE},        % [2] -> light blue
  morekeywords=[3]{SD, File},          % [3] -> bold orange
  morekeywords=[4]{open, exists},      % [4] -> orange
  %% The lines below add a nifty box around the code %%
  frame=single, % or shadow, none, double
  %rulesepcolor=\color{arduinoBlue},
}

% Define a pseudocode style
\lstdefinestyle{pseudocode}{
    basicstyle=\ttfamily\small,          % Monospaced font
    keywordstyle=\color{blue}\bfseries,  % Keywords in blue and bold
    keywords={if, else, while, end, print},
    keywordstyle=[2]\color{green!50!black}\bfseries,
    keywords=[2]{=, True},
    commentstyle=\color{green!70!black}, % Comments in green
    comment=[l]{//},                    % Single-line comments start with //
    stringstyle=\color{red},              % Strings in red
    string=[b]{''}                      % Strings enclosed in quotes
    numbers=left,                        % Line numbers on the left
    numberstyle=\ttfamily\color{gray},       % Style for line numbers
    stepnumber=1,                        % Line number in increments of 1
    numbersep=5pt,                       % Distance from the line numbers to the code
    backgroundcolor=\color{white},       % Background color
    frame=single,                        % Single frame around the code
    rulecolor=\color{black},             % Color of the frame
    tabsize=4,                           % Sets default tabsize to 4 spaces
    breaklines=true                     % Allows line breaking
}

% Set global styles for tcolorbox
\newtcolorbox{sticky}[2][]{
  colback=yellow!25,                     % Background color
  colframe=yellow!50,                    % Frame color
  rounded corners=all,                    % Rounded corners
  boxrule=1mm,                           % Border thickness
  title={\textcolor{blue}{#2}},          % Title color set 
  fonttitle=\bfseries\scriptsize,        % Title font style
  fontupper=\footnotesize,               % Content font size
  before skip=10pt,                      % Space before the box
  after skip=10pt,                       % Space after the box
  #1                                      % Additional options
}

\begin{document}
\tableofcontents

%\dominitoc % To create the mini table of contents

% \Huge Getting Going in Arduino
% \normalsize ~

% \noindent\rule{\textwidth}{1pt}\vspace{.5em}

\chapter{Hardware}
%\minitoc

\section{Arduino}

The Arduino platform is a flexible, relatively easy to learn, microprocessor capable of sensing, communication, and control.  There are many different versions of the Arduino, but we will focus on the two most common and fitting versions for this class.  These are shown in Figure~\ref{fig:arduinos}.

\begin{figure}[h]
    \centering
    \begin{subfigure}[b]{0.45\textwidth}
        \includegraphics[width=\textwidth]{figures/uno.jpg} % https://store-usa.arduino.cc/cdn/shop/files/A000066_03.front_1000x750.jpg?v=1727102657
        \caption{Arduino Uno R3, the most popular Arduino hardware.}
        \label{fig:uno}
    \end{subfigure}
    ~
    \begin{subfigure}[b]{0.45\textwidth}
        \includegraphics[width=\textwidth]{figures/nano.jpg} % https://store.arduino.cc/cdn/shop/files/A000005_03.front_1000x750.jpg?v=1753434497
        \caption{Arduino Nano, a much smaller format.} % (MKR instead?)
        \label{fig:nano}
    \end{subfigure}
    \caption{Arduino Hardware, which is constanly evolving.  These two versions are likely to be continue to be available in some format in the near future.}\label{fig:arduinos}
\end{figure}

You can think of the Arduino is a super simple computer, yet it's still different from the computer you probably use on a daily basis.  The Arduino has a set of general purpose input and output pins (GPIO).  This allows input and output of analog, digital, and communication protocols.  Not all pins support all function though!  See Figure~\ref{fig:arduino-pins} for details on the pins of the Uno.  Don't worry about all of the acronyms in there yet though, as we'll go through those in the next few pages here.

\begin{figure}
  \includegraphics[width=\textwidth]{figures/arduino-pins.png} % https://www.teachmemicro.com/wp-content/uploads/2017/03/arduino-pins.png
  \caption{The Arduino Uno platform has 19 GPIO pins available for use.  However, some of these pins support secondary functions.  For example, you can see that A4 and A5 on the lower left are labeled as SDA ADC4 and SCL ADC5 respectively.  SDA and SCL are communication type pins (discussed later).  ADC are analog to digital converter pins (also discussed later).  These pins cannot be used for both of these functions simultaneously. \label{fig:arduino-pins}}
\end{figure}

\chapter{Coding your first few Sketches}
%\minitoc % This will include sections of this chapter only

\section{Sketch Outline}
We call Arduino programs sketches. 
There are three major parts to Arduino code layout.  Let's start with the {\em second} section.

% \begin{lstlisting}[firstnumber=4,style=myArduino]
% void setup() {
%   // put your setup code here, to run once:

% }
% \end{lstlisting}

In this second section, as stated in the comments, code will run one time at start up.  This is where we will put our code that defines how pins on the General Purpose Input/Output (GPIO) will behave.  This could be setting a pin to be used as output, input, or even a communication interface.

In the {\em third} section we'll put the code that will run over and over again.  This comes right after the setup() loop.

% \begin{lstlisting}[firstnumber=4,style=myArduino]
% void setup() {
%   // put your setup code here, to run once:

% }

% void loop() {
%   // put your main code here, to run repeatedly:

% }
% \end{lstlisting}

Finally, we need to declare variables that we want to be scoped globally.  What does this mean?  Any variables that we want to be able to access in the setup() and the loop().  If you're new to programming an Arduino, then it might be best to just assume that all variables should be global for now.  There are indeed cases where you will want a variable to be accessible only inside of one of these, or a different function that you might write, but that's for later...

\begin{minipage}{\linewidth}
\begin{lstlisting}[style=myArduino]
// put your variable declarations here.
// and your includes.

void setup() {
  // put your setup code here, to run once:

}

void loop() {
  // put your main code here, to run repeatedly:

}
\end{lstlisting}
\end{minipage}

\section{Output}
\subsection{Blinking a Light Emitting Diode (LED)}
~
\marginnote{
  \begin{sticky}{What about Ohm's Law?}
  {We're going to make some assumptions for the moment, and will come back to this in~\ref{sec:ohms}.}
  \end{sticky}
}
To start with, we'll get an LED connected to the Arduino to blink.  Most Arduinos have a built in LED, but we're going to plug one in as a matter of practice.  As shown here, we are using a breadboard and a resistor, in addition to the LED.  However, in practice you can skip the breadboard and the resistor and connect your LED directly to the Arduino.

\begin{figure}[h]
    \centering
    \begin{subfigure}[b]{0.45\textwidth}
        \includegraphics[width=\textwidth]{figures/led-circuit.png} % https://docs.arduino.cc/built-in-examples/basics/Blink/
        \caption{Following best practice, we should have a resistor in our circuit.}
        \label{fig:led-circuit}
    \end{subfigure}
    ~
    \begin{subfigure}[b]{0.45\textwidth}
        \includegraphics[width=\textwidth]{figures/led-circuit-lazy.jpg} % https://www.makerguides.com/how-to-blink-an-led-using-arduino-4-different-ways/
        \caption{But we can also generally get away with being a little lazy and skipping the resistor.  More on this later!} % (MKR instead?)
        \label{fig:led-circuit-lazy}
    \end{subfigure}
    \caption{Arduino LED examples.}\label{fig:leds}
\end{figure}


Here's the code; let's look at what's really going on here.

\begin{minipage}{\linewidth}
\begin{lstlisting}[style=myArduino]
void setup() {
  pinMode(13, OUTPUT); // Set pin 13 as an output
}

void loop() {
  digitalWrite(13, HIGH); // Turn the LED on
  delay(1000); // Wait for one second
  digitalWrite(13, LOW); // Turn the LED off
  delay(1000); // Wait for one second
}
\end{lstlisting}
\end{minipage}

First up, we are not using any variables to store state or values.  We also are not using any special libraries, other supporting code, to do anything special.  Thus, line 1 starts with our setup() call.  Following this, we set the pin mode of GPIO pin 13 to be an output.

% \begin{lstlisting}[style=myArduino]
% void setup() {
%   pinMode(13, OUTPUT); // Set pin 13 as an output
% }
% \end{lstlisting}
\marginnote{
  \begin{sticky}{Is there a state between ``HIGH'' and ``LOW''?}
  Yes, but we need to save that for the section on pulse width modulation.
  \end{sticky}
}
Next up, we have the loop().  These four lines turn the LED on by setting GPIO13 to ``HIGH'', pause for 1 second (1000ms), turn the LED off by setting GPIO13 to ``LOW'', pause for another 1 second, and then repeats. When we set a pin to ``HIGH'' we are setting it to ``on'', or thinking in binary, ``TRUE'' at 5 volts (assuming we are using a 5 volt Arduino).  Likewise, setting the pin to ``LOW'', which you can also think of as ``FALSE'', sets the pin to ``off''.

% \begin{lstlisting}[firstnumber=5,style=myArduino]
% void loop() {
%   digitalWrite(13, HIGH); // Turn the LED on
%   delay(1000); // Wait for one second
%   digitalWrite(13, LOW); // Turn the LED off
%   delay(1000); // Wait for one second
% }
% \end{lstlisting}

\subsection{Adding a counter}
In this course, we are using the Arduino as a platform for collecting data.  Blinking an LED isn't really collecting data, or least not unless you want to sit by the LED and count how many times it blinks before the battery running the Arduino gives up.  So let's introduce another function that is built into the Arduino, one that will let us communicate with the Arduino.

For now, we have been plugging our Arduino into our computer to upload our programs, and incedentially, to power the Arduino.  With the introduction of this next function, we are assuming that is still the case.
\marginnote{
  \begin{sticky}{What is baud?}
  Baud, or baud rate, is how fast we will communicate between devices.  This is bits per second.
  \end{sticky}
}

\begin{minipage}{\linewidth}
\begin{lstlisting}[style=myArduino]
int ledBlinks = 0;     // Create a variable to store the number of times the LED has blinked for us

void setup() {
  pinMode(13, OUTPUT); // Set pin 13 as an output
  Serial.begin(9600);  // Define that we will be using the serial port, and 9600 baud
}

void loop() {
  digitalWrite(13, HIGH); // Turn the LED on
  delay(1000); // Wait for one second
  digitalWrite(13, LOW); // Turn the LED off
  delay(1000); // Wait for one second
  ledBlinks = ledBlinks + 1;
  Serial.print(ledBlinks);
}
\end{lstlisting}
\end{minipage}

\subsection{Using the Plotter}
The Arduino IDE has a built in data plotter.  This is especially helpful as we think about reading sensor data in realtime and debugging.  Let's modify our counter sketch to instead count up to some number, then count back down to another, then up again, and so on.  We are going to make a saw tooth function.

\begin{minipage}{\linewidth}
Saw Tooth Algorithm
% Use the pseudocode style without line numbers
\begin{lstlisting}[style=pseudocode]
x = 0
x_max = 20
x_min = 0
direction = 1
while True
  if x >= x_max
    direction = -1
  else if x <= x_min
    direction = 1
  end if
  x = x + direction
  print x
end while
\end{lstlisting}
\end{minipage}

This saw tooth algorithm should oscilate x between 0 and 20.  To translate this algorithm into Arduino code, we need to establish our variables (lines 1 so 5), setup the serial communication (lines ), and then loop through evaluating the state of x while making changes to the direction (lines ) and thus the value of x as appropriate (lines ).

\marginnote{
  \begin{sticky}{Ints}
  Arduino Ints are made up of two bytes and can range from -32,768 to 32,767.
  \end{sticky}
}

\begin{minipage}{\linewidth}
\begin{lstlisting}[style=myArduino]
int x = 0;     // to store the value we want to be outputting
int x_max = 20; // we can make this a larger number, try it after trying 20 first
int x_min = 0; // likewise to x_max, this can be changed after you give 0 a shot
int direction = 1; // see the note on ints

void setup() {
  // we don't really need the LED in this example, though you could use two LEDs to indicate the direction
  Serial.begin(9600);  // 9600 baud
}

void loop() {
  if x >= x_max {
    direction = -1;
  } else {
    direction = 1;
  }
  x = x + direction;
  Serial.print(x);
  delay(500); // Wait for half a second
}
\end{lstlisting}
\end{minipage}

Then show an image of the plotter in action.

\section{Input}
\subsection{Test for Echo}
take some input, and echo that back

\subsection{Calculator}
take two single digit inputs, and do addition

\subsection{Change Parameters}
take input and make things behave differently

restart and the defaults

\subsection{Sensors}
maybe start with a digital sensor, like an accelerometer?

then introduce analog in the next chapter?

\chapter{Analog Sensing}
Most all sensors are, at heart, analog in some way.  May of them, such as the accelerometer we used in the previous chapter, have an analog to digital converter built in.  But not all do.  Before we can discuss analog sensors in great detail, we need to cover some basic electrical engineering basics.

\section{Theory of Ohm's Law}
\label{sec:ohms}
overview of ohms law

\section{Ohm's Law Applied}
then into resistor potentiometers, and resistor networks
\subsection{Potentimeters}

\subsection{Resistor Networks}

with how to read resistors

\subsection{Variable Restistance Sensors}
such as flex sensors, force sensors, light sensors, etc

\chapter{Keeping Time}

\chapter{Saving the Data}

\end{document}
